\documentclass[11pt]{article}

\usepackage[margin=1in]{geometry}
\usepackage[T1]{fontenc}
\usepackage[utf8]{inputenc}
\usepackage{lmodern}
\usepackage{amsmath,amssymb,amsthm,bm,mathtools}
\usepackage{booktabs}
\usepackage{graphicx}
\usepackage{caption}
\usepackage{xcolor}
\usepackage{hyperref}
\usepackage{enumitem}
\usepackage{float}
\usepackage{algorithm}
\usepackage{algpseudocode}

\hypersetup{colorlinks=true, linkcolor=blue!60!black, citecolor=blue!60!black,
            urlcolor=blue!60!black}

\newtheorem{theorem}{Theorem}
\newtheorem{proposition}{Proposition}
\newtheorem{lemma}{Lemma}
\newtheorem{corollary}{Corollary}
\theoremstyle{remark}
\newtheorem{remark}{Remark}

\newcommand{\R}{\mathbb{R}}
\newcommand{\E}{\mathbb{E}}
\newcommand{\Var}{\operatorname{Var}}
\newcommand{\Cov}{\operatorname{Cov}}
\newcommand{\KL}{\operatorname{KL}}
\newcommand{\Normal}{\mathcal{N}}
\newcommand{\dd}{\,\mathrm{d}}
\newcommand{\Xibf}{\bm{\Xi}}

\title{Learning the Denoiser by Expectation Maximization\\
\large Belief propagation as an exact E-step for Markov diffusion priors}
\author{Thesis project --- Layer 5 working note}
\date{July 2026}

\begin{document}
\maketitle

\begin{abstract}
Earlier layers of this project established that when the clean data is a Markov
chain, belief propagation computes the diffusion score \emph{exactly}: the
posterior induced by coordinatewise noising is again a chain, chains are trees,
and sum--product BP on a tree is exact. That result assumed the prior was known.
This note removes the assumption. We treat the transition kernel of the clean
chain as an unknown parameter, observe only noised sequences, and estimate the
kernel by maximum marginal likelihood via expectation maximization.

Four things make this more than a routine application of EM. First, the E-step
is \emph{exact}, not variational: the same tree structure that made the score
exact makes the posterior expectations exact. Second, the entire E-step over an
arbitrarily large dataset compresses into a single $M \times M$ matrix
$\Xibf$, the continuous analogue of the expected transition-count matrix of
Baum--Welch, which is a sufficient statistic independent of how the kernel is
parameterized. Third, Fisher's identity turns $\Xibf$ into the exact gradient of
the marginal log-likelihood, so \emph{nothing is ever differentiated through the
forward--backward recursion} --- BP supplies the expectation rather than being
backpropagated through. Fourth, and this is the point of the exercise, the
learned object lives on $\R \times \R$ and carries no dependence on the
diffusion time: one fit yields the exact Bayes denoiser at \emph{every} noise
level, whereas a score network must learn a function on $\R^n \times \R_+$ and
be trained across the whole schedule.

We also make precise a remark that motivated the exercise --- that learning the
prior ``does not require noising at all''. Noising never destroys
identifiability, because the Gaussian characteristic function has no zeros
(Proposition~\ref{prop:ident}); what it destroys is \emph{information}, at a
rate we compute and measure. Section~\ref{sec:efficiency} states the efficiency
argument as a risk decomposition and Section~\ref{sec:limits} collects the
limitations, of which the honest one is that BP is dramatically slower than a
network at inference time.
\end{abstract}

\tableofcontents

% =====================================================================
\section{Setting}
\label{sec:setting}

\subsection{The two processes}

Throughout, the clean datum is a real sequence $a = (a_1,\dots,a_n) \in \R^n$
drawn from a prior $p_0$, and the forward (noising) process is the
coordinatewise variance-preserving Ornstein--Uhlenbeck SDE
\begin{equation}
  \dd X_s = -X_s \dd s + \sqrt{2}\dd W_s, \qquad X_0 = a,
\end{equation}
whose transition kernel at diffusion time $t$ is
\begin{equation}
  X_t \mid a \;\sim\; \Normal(\alpha_t a, \Delta_t I),
  \qquad
  \alpha_t = e^{-t}, \quad \Delta_t = 1 - e^{-2t}.
  \label{eq:ou}
\end{equation}
For unit-variance clean data $\alpha_t^2 + \Delta_t = 1$, so the noisy marginal
has unit variance at every $t$. We write $x$ for a realization of $X_t$ and
$p_t$ for its density.

Two facts from the earlier layers are used without reproof. The first is the
\emph{score identity}: because the channel \eqref{eq:ou} is linear-Gaussian,
\begin{equation}
  s(x,t) \;=\; \nabla_x \log p_t(x)
  \;=\; -\,\frac{x - \alpha_t\, m(x,t)}{\Delta_t},
  \qquad
  m(x,t) \;=\; \E[\,a \mid X_t = x\,],
  \label{eq:score-identity}
\end{equation}
so estimating the score and estimating the posterior mean (the
\emph{denoiser}) are the same problem, and an error in one transfers to the
other by the exact relation
$\hat s - s = (\alpha_t/\Delta_t)(\hat m - m)$.
The second is the \emph{exactness result}: if $p_0$ is Markov on the chain
$1 - 2 - \cdots - n$, then so is the posterior $p_t(a\mid x)$, and sum--product
BP returns its marginals exactly.

\subsection{The prior, now with unknown parameters}

We take the clean prior to be a chain with a parameterized transition kernel,
\begin{equation}
  p_\theta(a) \;=\; \mu(a_1) \prod_{i=2}^{n} K_\theta(a_i \mid a_{i-1}),
  \qquad \theta \in \Theta \subseteq \R^P,
  \label{eq:prior}
\end{equation}
with $\mu$ known. (Learning $\mu$ requires no new machinery --- it uses the
single-site belief at $i=1$ in place of the pairwise beliefs --- and is omitted
only to keep the notation light. In the experiments $\mu = \Normal(0,1)$ is
genuinely the law of $a_1$ under the data-generating process, so this is not an
approximation.)

The data available for learning is a sample of \emph{noisy} sequences,
\begin{equation}
  \mathcal{D} \;=\; \bigl\{\, (x^{(d)}, t_d) \,\bigr\}_{d=1}^{N},
  \qquad
  x^{(d)} \sim p_{t_d,\theta^\star}, \quad
  a^{(d)} \sim p_{\theta^\star} \text{ never observed}.
\end{equation}
Note what is and is not known: the observation channel \eqref{eq:ou} is known
exactly, including $t_d$, and the \emph{structure} of the prior (a chain) is
known; the \emph{content} of the prior, $K_\theta$, is not. This is the
situation the project cares about --- ``the data is Markov plus noise, but we do
not know the model behind it'' --- and it is strictly harder than the setting of
the earlier layers, where $K$ was handed to BP.

The estimator is maximum marginal likelihood,
\begin{equation}
  \hat\theta \;=\; \arg\max_{\theta}\; L(\theta),
  \qquad
  L(\theta) \;=\; \sum_{d=1}^{N} \log p_{t_d,\theta}\bigl(x^{(d)}\bigr),
  \qquad
  p_{t,\theta}(x) = \int_{\R^n} p_\theta(a) \prod_{i=1}^{n} \ell^t_i(a_i; x_i) \dd a,
  \label{eq:mml}
\end{equation}
where the per-site likelihood factor implied by \eqref{eq:ou} is
\begin{equation}
  \ell^t_i(a_i; x_i) \;=\; \Normal\bigl(x_i;\, \alpha_t a_i,\, \Delta_t\bigr).
  \label{eq:likelihood}
\end{equation}
The integral in \eqref{eq:mml} is over the unobserved clean chain: this is a
latent-variable problem, and EM is the natural way to attack it.

\begin{remark}[Why this is the right kind of ``learning'']
The hypothesis space here is not a function class but the set of Markov chains.
The estimator cannot represent a non-Markov prior at all, and it cannot
represent a denoiser that fails to be the exact Bayes denoiser of \emph{some}
chain. That is a severe restriction, and it is exactly the point: the
restriction is known to be satisfied (or nearly so) by the data, so imposing it
converts a nonparametric function-estimation problem into a low-dimensional
parametric one.
\end{remark}

% =====================================================================
\section{Exact EM through belief propagation}
\label{sec:em}

\subsection{The EM decomposition}

Fix a current parameter $\theta'$ and write $p_{\theta'}(a \mid x)$ for the
posterior over the clean chain. The standard decomposition of the marginal
log-likelihood reads, for a single observation $x$ at level $t$,
\begin{equation}
  \log p_{t,\theta}(x)
  \;=\;
  \underbrace{\E_{\theta'}\!\bigl[\log p_\theta(a,x)\bigr]}_{\textstyle Q(\theta\mid\theta')}
  \;+\;
  \underbrace{H\bigl(p_{\theta'}(\cdot\mid x)\bigr)}_{\text{free of }\theta}
  \;+\;
  \KL\bigl(p_{\theta'}(\cdot \mid x) \,\big\|\, p_{\theta}(\cdot \mid x)\bigr),
  \label{eq:em-decomp}
\end{equation}
where $\E_{\theta'}$ is taken under $p_{\theta'}(a \mid x)$ and
$p_\theta(a,x) = p_\theta(a)\prod_i \ell^t_i(a_i;x_i)$. Since the likelihood
factors \eqref{eq:likelihood} do not involve $\theta$, they contribute a
constant, and summing over the dataset gives
\begin{equation}
  Q(\theta \mid \theta')
  \;=\;
  \sum_{d=1}^{N} \sum_{i=2}^{n}
  \E_{\theta'}\!\left[\log K_\theta\bigl(a^{(d)}_i \bigm| a^{(d)}_{i-1}\bigr)\right]
  \;+\; \text{const}.
  \label{eq:Q}
\end{equation}
Two things follow immediately, and they organize everything below. First,
$Q$ depends on the posterior only through the \emph{pairwise} marginals of
adjacent sites. Second, because the $\KL$ term in \eqref{eq:em-decomp} is
non-negative and vanishes at $\theta = \theta'$,
\begin{equation}
  Q(\theta \mid \theta') \ge Q(\theta' \mid \theta')
  \;\;\Longrightarrow\;\;
  L(\theta) \ge L(\theta'),
  \label{eq:monotone}
\end{equation}
which is the monotone-ascent guarantee. We use \eqref{eq:monotone} as the
project's correctness test: an error anywhere in the forward--backward
recursion, the pairwise accumulation, or the M-step generically breaks it.

\subsection{The E-step is exact}

\begin{proposition}[Exact E-step]
\label{prop:exact-estep}
Let $p_\theta$ be the chain prior \eqref{eq:prior} and let the observation
likelihood factorize sitewise as $\prod_i \ell_i(a_i)$. Then the posterior
\begin{equation}
  p_\theta(a \mid x) \;\propto\;
  \mu(a_1)\,\ell_1(a_1) \prod_{i=2}^{n} K_\theta(a_i \mid a_{i-1})\,\ell_i(a_i)
\end{equation}
is a Markov chain on the same graph, and sum--product BP computes its
single-site and pairwise marginals exactly.
\end{proposition}

\begin{proof}
The likelihood factors are unary, so they reweight node potentials without
introducing edges: the factor graph of the posterior is the factor graph of the
prior with modified node potentials. That graph is a chain, hence a tree, and
sum--product BP is exact on trees. Exactness of the pairwise marginals on tree
edges is the standard consequence for the edge beliefs.
\end{proof}

This is the structural payoff of the whole project. In a generic latent-variable
model the E-step is intractable and must be approximated --- by loopy BP, by
expectation propagation, by a variational bound --- and the monotonicity
guarantee \eqref{eq:monotone} is lost along with the exactness. Here it is not
approximated at all. The only error is the numerical representation of the
continuous messages, which Layer 2 pinned at the $10^{-15}$ floor for the grids
used.

Concretely, writing $L_i$ and $R_i$ for the forward and backward messages
arriving at site $i$ (each \emph{excluding} the local likelihood), the
recursions and beliefs are
\begin{align}
  L_{i+1}(a') &= \int K_\theta(a' \mid a)\, L_i(a)\, \ell_i(a) \dd a,
  &
  R_{i-1}(a) &= \int K_\theta(a' \mid a)\, R_i(a')\, \ell_i(a') \dd a',
  \label{eq:messages}\\[2pt]
  b_i(a) &\propto L_i(a)\, \ell_i(a)\, R_i(a),
  &
  b_{i-1,i}(a, a') &\propto L_{i-1}(a)\,\ell_{i-1}(a)\, K_\theta(a' \mid a)\,
                       \ell_i(a')\, R_i(a').
  \label{eq:beliefs}
\end{align}

\subsection{The whole E-step is one matrix}

Represent the messages on a grid $u_1,\dots,u_M$ with quadrature weights
$w_1,\dots,w_M$, and define for each chain $d$ and each site $i$ the vectors
\begin{equation}
  f^{(d)}_i(j) \;=\; w_j\, L^{(d)}_i(u_j)\, \ell^{(d)}_i(u_j),
  \qquad
  g^{(d)}_i(k) \;=\; w_k\, R^{(d)}_i(u_k)\, \ell^{(d)}_i(u_k),
\end{equation}
i.e.\ the incoming messages \emph{with} their local factors folded in. Writing
$K[k,j] = K_\theta(u_k \mid u_j)$, the pairwise belief \eqref{eq:beliefs} on
edge $(i,i+1)$ evaluated at the grid pair $(u_j, u_k)$ is
$f^{(d)}_i(j)\, K[k,j]\, g^{(d)}_{i+1}(k) / Z^{(d)}_i$ with
$Z^{(d)}_i = (g^{(d)}_{i+1})^{\!\top} K f^{(d)}_i$.

\begin{proposition}[Sufficiency of the pair-count matrix]
\label{prop:xi}
Define
\begin{equation}
  C \;=\; \sum_{d=1}^{N} \sum_{i=1}^{n-1}
          \frac{g^{(d)}_{i+1} \bigl(f^{(d)}_i\bigr)^{\!\top}}{Z^{(d)}_i}
  \;\in\; \R^{M \times M},
  \qquad
  \Xibf \;=\; C \odot K,
  \label{eq:xi}
\end{equation}
with $\odot$ the entrywise product. Then $\Xibf[k,j]$ is the total posterior
mass the dataset places on the transition $u_j \mapsto u_k$; it satisfies
$\sum_{j,k} \Xibf[k,j] = N(n-1)$; and the expected complete-data
log-likelihood \eqref{eq:Q} is the bilinear form
\begin{equation}
  Q(\theta \mid \theta') \;=\; \bigl\langle \Xibf(\theta'),\; \log K_\theta \bigr\rangle
  \;+\; \text{const},
  \qquad
  \langle A, B\rangle = \textstyle\sum_{j,k} A[k,j]\,B[k,j].
  \label{eq:Q-bilinear}
\end{equation}
\end{proposition}

\begin{proof}
Immediate from \eqref{eq:Q} and the grid form of \eqref{eq:beliefs}: each edge
contributes its normalized pairwise belief, which sums to one over the grid,
and $Q$ is the sum over edges of $\sum_{j,k} b_{i,i+1}(u_j,u_k)\log K_\theta[k,j]$.
Collecting the belief mass across edges and chains gives $\Xibf$.
\end{proof}

Proposition~\ref{prop:xi} is the algorithmic heart of the method, and it is
worth stating what it buys.

\begin{enumerate}[leftmargin=*]
\item \textbf{$\Xibf$ is a sufficient statistic of the E-step.} It does not
  depend on how $K_\theta$ is parameterized. A scalar AR(1), a mixture, a neural
  network --- all consume the same $M \times M$ matrix. Changing the
  parameterization changes only the M-step.
\item \textbf{The M-step never touches the data.} After the E-step, the dataset
  can be discarded. The M-step cost is $O(P M^2)$ regardless of $N$.
\item \textbf{Different noise levels simply add.} The level $t$ enters only
  through $\ell^t_i$; statistics from observations at different $t$ accumulate
  into one $\Xibf$. A dataset spread over the whole diffusion schedule is
  summarized by a single matrix.
\item \textbf{It is the continuum Baum--Welch statistic.} For a discrete-state
  HMM, $\Xibf$ is exactly the matrix of expected transition counts. The
  continuous case differs only in that the state space must be represented.
\end{enumerate}

\subsection{Gradients without differentiating through BP}

The M-step maximizes \eqref{eq:Q-bilinear}. When no closed form is available one
follows the gradient, and here the relevant gradient is available exactly and
cheaply.

\begin{proposition}[Fisher's identity for the chain]
\label{prop:fisher}
For any $\theta$ at which $\log K_\theta$ is differentiable,
\begin{equation}
  \nabla_\theta L(\theta)
  \;=\;
  \bigl\langle \Xibf(\theta),\; \nabla_\theta \log K_\theta \bigr\rangle,
  \label{eq:fisher}
\end{equation}
where $\Xibf(\theta)$ is the pair-count matrix \eqref{eq:xi} assembled at the
same $\theta$. In particular the gradient of the \emph{exact marginal
log-likelihood} requires one forward--backward pass and no differentiation of
the BP recursion.
\end{proposition}

\begin{proof}
For one observation,
$\nabla_\theta \log p_\theta(x)
 = p_\theta(x)^{-1}\!\int \nabla_\theta p_\theta(a,x)\dd a
 = \int p_\theta(a\mid x)\, \nabla_\theta \log p_\theta(a,x) \dd a
 = \E_{\theta}\bigl[\nabla_\theta \log p_\theta(a)\bigr]$,
the likelihood factors being $\theta$-free. Since
$\nabla_\theta \log p_\theta(a) = \sum_i \nabla_\theta \log K_\theta(a_i\mid a_{i-1})$
involves only adjacent pairs, the expectation is taken under the pairwise
posteriors, and summing over the dataset gives \eqref{eq:fisher}. Equivalently:
$\nabla_\theta L(\theta) = \nabla_\theta Q(\theta \mid \theta')|_{\theta'=\theta}$.
\end{proof}

\begin{remark}[This is the computational point]
It is tempting to implement ``learning inside BP'' by writing the
forward--backward recursion in an autodiff framework and backpropagating the
evidence through it. Proposition~\ref{prop:fisher} says this is unnecessary and
wasteful: the derivative of the log-evidence with respect to the parameters is
an \emph{expectation under the beliefs BP already computed}. BP is not a
computation graph to differentiate; it is the oracle that supplies the
expectation. The implementation accompanying this note is plain \texttt{numpy},
with no autodiff anywhere, and its gradients agree with finite differences of
the exact evidence to $\sim\!10^{-9}$ relative for smooth kernels
(Section~\ref{sec:kernels} records the one exception).
\end{remark}

\subsection{The algorithm and its cost}

\begin{algorithm}[H]
\caption{EM for a chain prior observed through the OU channel}
\begin{algorithmic}[1]
\State \textbf{input} noisy chains grouped by level $\{(X^{(g)}, \alpha_g, \Delta_g)\}$,
       grid $u$, weights $w$, initial $\theta^{(0)}$
\For{$r = 0, 1, 2, \dots$}
  \State $K \gets \exp \log K_{\theta^{(r)}}(u)$
        \Comment{$O(P M^2)$, once per iteration}
  \State $\Xibf \gets 0$,\; $L \gets 0$
  \For{each group $g$ and each chain $d$ in it}
    \State forward--backward pass \eqref{eq:messages}
           \Comment{$O(n M^2)$ per chain, batched as matmuls}
    \State accumulate $\Xibf$ by \eqref{eq:xi} and the exact
           $\log p_{t_g,\theta^{(r)}}(x^{(d)})$ into $L$
  \EndFor
  \State $\theta^{(r+1)} \gets \arg\max_\theta \langle \Xibf, \log K_\theta\rangle$
         \Comment{closed form, or an ascent step (generalized EM)}
  \State \textbf{assert} $L$ non-decreasing
         \Comment{the correctness test \eqref{eq:monotone}}
\EndFor
\end{algorithmic}
\end{algorithm}

Per iteration the cost is $O(N n M^2)$ for the E-step and $O(P M^2)$ for the
M-step. The $M^2$ is the price of representing continuous messages; the
important scalings are that the cost is \emph{linear} in the dataset size and
\emph{independent of the parameter count} in the dominant term. A score network
inverts both: its cost per gradient step is proportional to the parameter count
and it needs many passes over the data.

% =====================================================================
\section{What noising costs, and what it does not}
\label{sec:noising}

The framing that motivated this note included the remark that learning the prior
parameters ``does not require noising at all''. That is correct, and it is worth
saying precisely why, because the precise version is more useful than the slogan.

\subsection{The two limits}

At $t \to 0$ we have $\alpha_t \to 1$, $\Delta_t \to 0$, and the likelihood
factor \eqref{eq:likelihood} tends to $\delta(a_i - x_i)$. The posterior
collapses onto the observation, the pairwise beliefs become products of deltas,
$\Xibf$ degenerates into a histogram of \emph{observed} transitions, and
\eqref{eq:Q-bilinear} becomes the ordinary complete-data log-likelihood. EM
degenerates to plain maximum likelihood on clean data --- no latent variables,
no iteration required for kernels with closed-form MLEs. At the other end,
$t \to \infty$ gives $\alpha_t \to 0$: the observation is independent of $a$,
the posterior equals the prior, and $\nabla_\theta L \to 0$. Nothing can be
learned.

So noising is not needed to identify the prior; the clean chain, if available,
identifies it more efficiently. Noising is what the \emph{problem} imposes when
only noisy data exists --- and, more interestingly, it is a knob we can turn.

\subsection{Identifiability survives any finite noise level}

\begin{proposition}[Noising preserves identifiability]
\label{prop:ident}
Fix $t \in (0,\infty)$ and let $P_\theta$ denote the law of the clean sequence
$a \in \R^n$ under \eqref{eq:prior}. If $p_{t,\theta} = p_{t,\theta'}$ as
distributions on $\R^n$, then $P_\theta = P_{\theta'}$. Consequently the model
is identifiable at noise level $t$ if and only if it is identifiable from clean
data.
\end{proposition}

\begin{proof}
Under \eqref{eq:ou}, $X_t = \alpha_t a + \sqrt{\Delta_t}\, z$ with
$z \sim \Normal(0,I)$ independent of $a$, so the characteristic functions
satisfy
\begin{equation}
  \varphi_{t,\theta}(\xi)
  \;=\; \varphi_\theta(\alpha_t \xi)\, e^{-\Delta_t \|\xi\|^2 / 2},
  \qquad \xi \in \R^n,
\end{equation}
where $\varphi_\theta$ is the characteristic function of $P_\theta$. The
Gaussian factor $e^{-\Delta_t\|\xi\|^2/2}$ is strictly positive for every
$\xi$, and $\alpha_t = e^{-t} > 0$ for every finite $t$. Hence
$\varphi_{t,\theta} = \varphi_{t,\theta'}$ forces
$\varphi_\theta(\alpha_t\xi) = \varphi_{\theta'}(\alpha_t\xi)$ for all $\xi$,
i.e.\ $\varphi_\theta = \varphi_{\theta'}$ on $\R^n$, and characteristic
functions determine distributions.
\end{proof}

\begin{remark}[Injective but ill-conditioned]
Proposition~\ref{prop:ident} is a statement about population distributions, and
it is exactly as strong --- and as weak --- as the classical fact that Gaussian
deconvolution is injective. Injectivity does not mean stability: the inverse map
is unbounded, and dividing by $e^{-\Delta_t\|\xi\|^2/2}$ amplifies high
frequencies catastrophically. Finite-sample behaviour is therefore governed not
by identifiability but by \emph{information}, which is the subject of the next
subsection. The practical reading is: no noise level is fatal in principle, but
the estimator's variance grows quickly with $t$, and features of $P_\theta$
living at high frequency --- sharp peaks, heavy tails, the very non-Gaussianity
this project is about --- are the first to be lost.
\end{remark}

\subsection{The information budget}

The relevant finite-sample quantity is the Fisher information of $\theta$
carried by one noisy chain,
\begin{equation}
  J_t \;=\; \E_{x \sim p_{t,\theta}}\!\left[
    \nabla_\theta \log p_{t,\theta}(x)\,
    \nabla_\theta \log p_{t,\theta}(x)^{\!\top}\right],
  \label{eq:fisher-info}
\end{equation}
which by Proposition~\ref{prop:fisher} is computable exactly from single-chain BP
passes: the per-chain score is $\langle \Xibf(x), \nabla_\theta \log K_\theta\rangle$,
requiring no simulation of the estimator and no numerical differentiation. The
Cram\'er--Rao bound then gives $\Var(\hat\theta_p) \gtrsim [J_t^{-1}]_{pp}/N$
for unbiased estimators, which we use as an order-of-magnitude yardstick rather
than a strict bound (the maximum-likelihood estimator is biased at finite $N$,
and the measured bias is reported alongside the variance in the experiments).

$J_t$ is monotonically destroyed as $t$ grows. The measured decay, together with
the realized estimator error, is the content of Part~2 of experiment~06:
the point of that comparison is to check the estimator against \emph{its own
information budget}, so that a large error at large $t$ is correctly read as
``the data does not contain this'' rather than ``the algorithm is failing''.

\begin{corollary}[The remark, made precise]
If clean sequences are available, estimate $\theta$ from them directly: this is
the $t\to0$ limit, it is the most informative case, and it requires no BP at all
during fitting. The exact denoiser at \emph{every} noise level then follows by
running BP with $\hat\theta$, with no further learning. If only noisy sequences
are available, the same estimator is still consistent for any finite $t$
(Proposition~\ref{prop:ident}) and EM with an exact BP E-step computes it.
\end{corollary}

Both regimes are implemented and compared: \texttt{fit\_clean} builds $\Xibf$ as
a histogram of observed transitions, \texttt{fit\_em} builds it from beliefs, and
the downstream M-step code is identical.

% =====================================================================
\section{Parameterizations and their M-steps}
\label{sec:kernels}

Proposition~\ref{prop:xi} isolates all the modelling choice into the M-step. We
implement four rungs of expressivity, in increasing order.

\subsection{Rung 1 --- Gaussian AR(1): the correctness harness}

$K_\theta(a'\mid a) = \Normal(a'; \rho a, q)$, $\theta = (\rho, q)$. Setting
$\nabla_\theta Q = 0$ in \eqref{eq:Q-bilinear} gives, with
$S_{ab} = \sum_{j,k}\Xibf[k,j]\, u_j^{a} u_k^{b}$ and $W = \sum_{j,k}\Xibf[k,j]$,
\begin{equation}
  \rho^{+} = \frac{S_{11}}{S_{20}},
  \qquad
  q^{+} = \frac{S_{02} - (S_{11})^2/S_{20}}{W},
  \label{eq:gauss-mstep}
\end{equation}
the belief-weighted Yule--Walker equations. In the $t\to0$ limit $\Xibf$ becomes
a transition histogram and \eqref{eq:gauss-mstep} reduces to ordinary AR(1)
least squares, as it must. This rung exists to be checked against closed forms:
the whole prior is Gaussian, so the evidence, the score and the posterior mean
are all available analytically.

\subsection{Rung 2 --- Laplace AR(1): closed form despite a non-smooth kernel}

$K_\theta(a'\mid a) = \tfrac{1}{2b}\exp(-|a' - \rho a|/b)$, $\theta = (\rho,b)$.
Here
\begin{equation}
  Q(\rho, b) \;=\; -\frac{1}{b}\sum_{j,k}\Xibf[k,j]\,\bigl|u_k - \rho u_j\bigr|
                  \;-\; W\log 2b .
\end{equation}
For fixed $\rho$ the optimal scale is the belief-weighted mean absolute
residual, $b(\rho) = W^{-1}\sum \Xibf[k,j]|u_k - \rho u_j|$; substituting back
gives $Q(\rho) = -W\log 2b(\rho) - W$, so maximizing over $\rho$ is
\emph{minimizing a weighted mean absolute residual} --- a least-absolute-
deviations regression through the origin. Its exact solution is the weighted
median of the ratios $u_k/u_j$ with weights $\Xibf[k,j]\,|u_j|$. The joint
M-step is therefore closed form even though the kernel is not differentiable:
no line search, no gradient.

\begin{remark}[A discretization caveat, and why it is not generic]
\label{rem:quantization}
Two things go wrong for this kernel on a uniform grid through the origin, and
both are artifacts of the discretization rather than of EM.

\emph{(i) Gradient accuracy.} $\partial_\rho \log K_\theta$ contains
$\operatorname{sign}(a' - \rho a)$, a discontinuity across the diagonal.
Trapezoidal quadrature of a discontinuous integrand loses the spectral accuracy
the rest of this package enjoys: against a finite difference of the exact grid
evidence, the $\rho$-component of \eqref{eq:fisher} is off by $\sim\!13\%$ at
$M = 201$ and still $\sim\!0.1\%$ at $M = 1601$, decaying erratically. The
$b$-component, which is smooth, agrees to machine precision --- confirming that
Proposition~\ref{prop:fisher} is fine and the quadrature is not.

\emph{(ii) Lattice snapping.} The minimizer of a weighted sum of
$|u_k - \rho u_j|$ is a breakpoint, i.e.\ one of the ratios $u_k/u_j$. On a
uniform grid through the origin these ratios are rationals $m/l$, and a
low-denominator value such as $4/5$ has many aliases ($8/10$, $12/15$, \dots)
that pool weight onto it. So $\hat\rho$ lands on simple rationals: spuriously
\emph{exactly} correct when the truth happens to be one, and stuck at lattice
resolution when it is not. Experiment~06 Part~5 measures this against $M$ and
against $\rho^\star$; the reported $N^{-1/2}$ rates use the smooth kernels,
whose M-steps are ratios of smooth moments and vary continuously.
\end{remark}

\begin{remark}[The literal gradient route, measured]
\label{rem:gradient-route}
The suggestion that started this layer was phrased as gradient ascent --- iterate
along $\nabla_\theta L$ from a random start --- rather than as EM. Fisher's
identity makes that route immediately available, since \eqref{eq:fisher} is the
gradient of the exact marginal likelihood from the same single BP pass. Both
routes share the E-step and the statistic $\Xibf$; they differ only in what they
do with it:
\begin{equation}
  \text{gradient:}\quad \theta \leftarrow \theta + \eta\,\langle\Xibf, \nabla_\theta\log K_\theta\rangle,
  \qquad
  \text{EM:}\quad \theta \leftarrow \arg\max_\theta \langle\Xibf, \log K_\theta\rangle .
\end{equation}
Experiment~08 measures them side by side ($256$ chains, $n = 24$, $120$
iterations), and the answer is not one-sided.

On the smooth Gaussian kernel the gradient route \emph{converges to the same
place}: at $\eta = 0.5$ it matches EM's final log-likelihood to $2\times10^{-9}$
nats and its parameter errors to four digits. At $\eta = 0.1$ it is simply not
there yet ($1.6$ nats short); at $\eta = 2$ it diverges, violating monotonicity
by $527$ nats; at $\eta = 8$ it leaves the admissible set entirely. Nothing is
gained over EM here, and a step size must be tuned to lose nothing.

On the Laplace kernel the gradient route is genuinely \emph{better}. It reaches
a higher log-likelihood than the exact M-step ($-8064.7$ against $-8066.6$ at
$\eta = 0.1$) and better parameters on both coordinates ($|\Delta\rho| = 0.022$
against $0.050$, $|\Delta b| = 0.013$ against $0.027$). The reason is
Remark~\ref{rem:quantization}: EM's M-step here is exactly optimal for the
discretized model and lands on a lattice point, reaching a genuine fixed point
after $72$ iterations with $\hat\rho$ stuck $0.05$ from the truth, while
gradient steps move continuously and walk past it. Divergence, when it happens,
is unmistakable rather than subtle --- at $\eta = 8$ the scale collapses toward
a near-delta kernel and the reported log-density goes positive.

So the two are complementary rather than ranked, and the split is not the one
intuition suggests. The exact M-step is step-size-free and exactly monotone, and
should be preferred wherever it exists \emph{and} the kernel is smooth. The
gradient route earns its place precisely where the M-step is least trustworthy:
where it has no closed form (Rung 4), or where its exact solution is an artifact
of the discretization rather than of the model (Rung 2). What is common to both,
and is the point of Proposition~\ref{prop:fisher}, is that neither ever
differentiates through the forward--backward recursion.
\end{remark}

\subsection{Rung 3 --- unknown innovation law}

Keep the linear autoregression but make the innovation density essentially free:
\begin{equation}
  K_\theta(a' \mid a) \;=\; \sum_{c=1}^{C} \pi_c\,
     \Normal\bigl(a' - \rho a;\, \mu_c,\, \sigma_c^2\bigr),
  \qquad
  \theta = \bigl(\rho, \{\pi_c,\mu_c,\sigma^2_c\}_{c=1}^C\bigr).
\end{equation}
This is the estimator to use when the model behind the data is genuinely
unknown --- it has never heard of the Laplace density it will be asked to
recover. The complete-data problem now carries a \emph{second} latent variable,
the mixture label, so the M-step is itself an EM, run against the fixed $\Xibf$.
With responsibilities $r_c(j,k)$ computed at the current parameters and
$e_{jk} = u_k - \rho u_j$, each block is closed form:
\begin{equation}
  \pi_c^{+} = \frac{\langle \Xibf, r_c\rangle}{W},
  \quad
  \mu_c^{+} = \frac{\langle \Xibf, r_c e\rangle}{\langle \Xibf, r_c\rangle},
  \quad
  (\sigma^2_c)^{+} = \frac{\langle \Xibf, r_c (e - \mu^+_c)^2\rangle}{\langle \Xibf, r_c\rangle},
  \quad
  \rho^{+} = \frac{\sum_c \sigma_c^{-2}\langle \Xibf, r_c\, u_j (u_k - \mu_c)\rangle}
                  {\sum_c \sigma_c^{-2}\langle \Xibf, r_c\, u_j^2 \rangle}.
  \label{eq:mix-mstep}
\end{equation}
Alternating these blocks is expectation conditional maximization, which is
monotone; since $\Xibf$ is fixed, each inner sweep is $O(CM^2)$ on a matrix that
does not grow with the dataset.

\begin{remark}[A constraint that looks free and is not]
The innovation of a zero-mean stationary chain should satisfy
$\sum_c \pi_c \mu_c = 0$. Imposing it by recentering the component means after
each update is \emph{not} harmless: the projection does not maximize $Q$, and in
testing it broke monotone ascent by $\sim 10^{-2}$ nats. Left unconstrained, ECM
is exactly monotone and the fitted innovation mean comes out at the $10^{-3}$
level by itself. We therefore report the innovation mean as a diagnostic rather
than enforcing it --- a small illustration of the general rule that every step
of an EM must be justified as an ascent step, not merely as a reasonable-looking
adjustment.
\end{remark}

\subsection{Rung 4 --- a neural transition kernel}

The final rung answers the ``small deep network'' suggestion directly, but
places the network \emph{inside} the exact recursion rather than in place of it.
Let a small MLP $\phi$ map a state to the parameters of a conditional mixture,
\begin{equation}
  K_\phi(a' \mid a) \;=\; \sum_{c=1}^{C} \pi_c(a)\,
    \Normal\bigl(a';\, m_c(a),\, \sigma^2_c(a)\bigr),
  \qquad
  \bigl(\pi, m, \log\sigma^2\bigr)(a) = \mathrm{MLP}_\phi(a),
\end{equation}
a mixture-density network, nonparametric in both the location and the shape of
the transition. There is no closed-form M-step, so we take one Adam step on $Q$
per E-step --- generalized EM, which only requires $Q$ to increase, and which
therefore requires \emph{checking} that it does; see
Remark~\ref{rem:backtrack}. The gradient is where the structure pays off. By the
chain rule applied to \eqref{eq:Q-bilinear},
\begin{equation}
  \frac{\partial Q}{\partial \phi}
  \;=\;
  \sum_{j} \left[
    \sum_{k} \widetilde{\Xibf}[k,j]\,
      \frac{\partial \log K_\phi[k,j]}{\partial\, \mathrm{head}(u_j)}
  \right]
  \frac{\partial\, \mathrm{head}(u_j)}{\partial \phi},
  \label{eq:mdn-grad}
\end{equation}
so \textbf{the network is evaluated and backpropagated only at the $M$ grid
points}, once per EM iteration, \emph{no matter how large the dataset is}. All
dependence on the data has been absorbed into $\Xibf$. A vanilla score network,
by contrast, must backpropagate through every training example at every step.
Here $\widetilde{\Xibf}[k,j] = \Xibf[k,j] - \bigl(\sum_k \Xibf[k,j]\bigr) p_j(k)$
is the centered statistic, $p_j$ being the grid-normalized kernel column; the
subtracted term is the derivative of the column log-partition, the familiar
``observed minus expected'' form. Verified against finite differences of the
exact evidence to $1.6 \times 10^{-6}$ relative.

\begin{remark}[A fixed learning rate is not an ascent step]
\label{rem:backtrack}
``Generalized EM only needs $Q$ to increase'' is a licence to take a cheap
M-step, not a licence to skip the check. Run with a fixed Adam step, this kernel
raised the marginal log-likelihood overall ($-18978 \to -16652$ nats on a
400-chain fit) while violating monotonicity by $547$ nats at $\eta = 0.02$ and
$1452$ at $\eta = 0.05$ --- at which point EM's one guarantee is gone and a
diverging run looks exactly like a slow one. Backtracking the step until $Q$
actually increases restores exact monotonicity ($0$ violation at both rates) at
negligible cost, since $Q = \langle\Xibf, \log K\rangle$ is one contraction of a
matrix already in hand. It is also diagnostic: at $\eta = 0.05$ the backtracked
run stalls at $-17392$ instead of silently descending, which is the learning
rate reporting itself as too large.
\end{remark}

\begin{remark}[More expressive is not better here]
\label{rem:rung4-loses}
On the Laplace chain with $400$ observed noisy chains ($M = 201$, $250$ EM
iterations), the neural kernel with $873$ parameters is dominated on
\emph{both} axes by the $13$-parameter mixture kernel of Rung 3: training
log-likelihood $-16781$ against $-16656$, and held-out relative score error
$0.052$ against $0.026$. It has not bought expressivity that this data needs; it
has bought optimization difficulty, since a gradient M-step under a monotonicity
safeguard makes slower progress per iteration than a closed-form one.

The one place the network does edge ahead is instructive and not a
recommendation: run \emph{without} the safeguard of
Remark~\ref{rem:backtrack} at $\eta = 0.05$, it attains $-16652$, marginally
above the mixture's likelihood, while its held-out error stays worse at
$0.059$ --- ordinary overfitting, reached by a procedure that violated
monotonicity by $1452$ nats to get there.

Rung 4 is the right tool when the transition genuinely has state-dependent
shape that Rung 3 cannot express. On data where a linear autoregression with a
free innovation law is correct, it costs two orders of magnitude in parameters
and returns nothing. The efficiency argument of Section~\ref{sec:efficiency} is
an argument for matching the hypothesis class to the structure --- and it
applies to the structured method's own internals, not only to the baseline it
is being compared against.
\end{remark}

% =====================================================================
\section{Why this should beat a score network}
\label{sec:efficiency}

\subsection{The decoupling}

The structural claim is one sentence: \emph{the learned parameters live on the
clean chain, and the noise level lives in the likelihood}. The object being
estimated, $K_\theta$, is a function on $\R \times \R$ with no dependence on $t$.
Once $\hat\theta$ is fixed, the denoiser at any noise level is obtained by
running BP with the likelihood factors \eqref{eq:likelihood} at that level ---
exactly, and with no further fitting. The estimation problem and the diffusion
schedule are decoupled.

A denoising-score-matching network has no such decoupling. It represents
$(x,t) \mapsto m$ as a function on $\R^n \times \R_+$; it must be trained across
the whole schedule at once; and the target it is fitting changes with $t$.
Nothing in the architecture knows that the chain is Markov, so the conditional
independencies must be recovered from data.

\subsection{Risk decomposition}

Write $m_\theta(x,t)$ for the exact posterior mean under prior $p_\theta$. For
the BP estimator the excess risk at level $t$ decomposes as
\begin{equation}
  \underbrace{\bigl\| m_{\hat\theta}(\cdot,t) - m_{\theta^\star}(\cdot,t)\bigr\|}_{\text{total}}
  \;\le\;
  \underbrace{\bigl\| m_{\hat\theta} - m_{\theta^\dagger}\bigr\|}_{\text{estimation}}
  \;+\;
  \underbrace{\bigl\| m_{\theta^\dagger} - m_{\theta^\star}\bigr\|}_{\text{family misspecification}},
  \label{eq:risk}
\end{equation}
where $\theta^\dagger$ is the population optimum within the chosen family.
Crucially there is \emph{no third term}: given a prior, BP returns the exact
Bayes denoiser, so there is no approximation error from a function class. The
misspecification term is zero when the family contains the truth (Rungs 1--2 on
matching data) and is the observed error floor otherwise (Rung 3 on Laplace
data, where it decreases with $C$). The estimation term is parametric.

\begin{proposition}[Sensitivity of the BP denoiser; a delta method]
\label{prop:sensitivity}
Suppose $\theta \mapsto \log K_\theta$ is $C^1$ and that, for
$t$ in a compact interval $[t_{\min}, t_{\max}]$ with $t_{\min} > 0$, the
posterior variances are bounded. Then $m_\theta(x,t)$ is differentiable in
$\theta$ with
\begin{equation}
  \frac{\partial m_{\theta,i}(x,t)}{\partial \theta_p}
  \;=\;
  \Cov_{p_\theta(\cdot\mid x)}\!\left(
     a_i,\;
     \sum_{l=2}^{n} \frac{\partial \log K_\theta(a_l \mid a_{l-1})}{\partial \theta_p}
  \right),
  \label{eq:sensitivity}
\end{equation}
a posterior covariance computable from the same beliefs BP already produces.
Consequently, if $\hat\theta - \theta^\dagger = O_P(N^{-1/2})$, then
$\sup_{t \in [t_{\min},t_{\max}]}
   \|m_{\hat\theta}(\cdot,t) - m_{\theta^\dagger}(\cdot,t)\| = O_P(N^{-1/2})$.
\end{proposition}

\begin{proof}[Proof sketch]
Differentiate $m_{\theta,i} = \int a_i\, p_\theta(a\mid x)\dd a$ under the
integral. Since
$\partial_{\theta_p} p_\theta(a \mid x)
 = p_\theta(a\mid x)\bigl[\partial_{\theta_p}\log p_\theta(a)
   - \E_\theta \partial_{\theta_p}\log p_\theta(a)\bigr]$
and $\log p_\theta(a)$ is the sum of log-transitions, \eqref{eq:sensitivity}
follows. The uniformity statement is the mean value theorem applied along the
segment from $\theta^\dagger$ to $\hat\theta$, using that the covariance
\eqref{eq:sensitivity} is bounded on the compact $t$-interval.
\end{proof}

The content of Proposition~\ref{prop:sensitivity} is the word \emph{uniformly}:
one $N^{-1/2}$-consistent parameter estimate delivers an $N^{-1/2}$-accurate
denoiser at \emph{every} noise level simultaneously. The network gets no such
statement; its error must be paid for separately in every region of $(x,t)$ it
is asked about, and in particular it must extrapolate when asked about a $t$
outside its training schedule.

\subsection{What the comparison must control}

An efficiency claim is only as good as the baseline, so the experiments are set
up to favour the network wherever there was a choice:
\begin{enumerate}[leftmargin=*]
\item \textbf{Data budget} is the number of \emph{clean chains} $N$. The network
  trains on paired $(a,x)$ with a fresh noise draw at every gradient step, so it
  may consume unlimited noise realizations of those $N$ chains. EM sees one
  noisy realization of each chain, at one level, and never sees a clean chain.
  The two methods are therefore \emph{not} on equal information: EM is given
  strictly less.
\item \textbf{Objective} is denoising score matching, whose minimizer is exactly
  $\E[a\mid x,t]$ --- the same target BP computes. The network is not
  handicapped by a surrogate loss. Both standard parameterizations are trained
  and reported: noise prediction ($\|z_\phi - z\|^2$, the usual diffusion
  recipe) and clean-signal prediction ($\|a_\phi - a\|^2$). They share a
  minimizer but not a finite-sample loss, the two differing by the weight
  $\alpha_t^2/\Delta_t$; noise prediction additionally recovers the mean as
  $(x - \sqrt{\Delta_t}\, z_\phi)/\alpha_t$, whose $\sqrt{\Delta_t}$ prefactor
  suppresses network error exactly where signal prediction is weakest. Resting
  the comparison on whichever parameterization happens to lose would not be
  evidence of anything.
\item \textbf{Capacity and training budget} are swept, so the comparison cannot
  hinge on an undersized or undertrained network.
\item \textbf{Evaluation} is on held-out chains against exact grid BP under the
  true prior, with the score/mean identity
  $\hat s - s = (\alpha_t/\Delta_t)(\hat m - m)$ verified to machine precision
  on every row as an implementation guard.
\end{enumerate}

% =====================================================================
\section{Limitations}
\label{sec:limits}

\begin{enumerate}[leftmargin=*]
\item \textbf{Inference cost is the honest loss.} Grid BP costs $O(nM^2)$ per
  chain per evaluation; a network forward pass is a few matmuls. Since reverse
  diffusion evaluates the denoiser at every integration step, this is not a
  detail. The measured slowdown is two orders of magnitude
  (experiment~07 Part~4) and is reported as prominently as the wins. The natural
  responses --- distilling the fitted BP denoiser into a network, or using the
  Gaussian-projected BP of Layer 2 at inference --- are the obvious next step,
  not something claimed here.
\item \textbf{Everything rests on the chain assumption.} If the prior is not
  Markov, the estimator is misspecified in a way no amount of data fixes.
  Layer~4 showed the correction is exactly rank one for AR(1)-plus-global-latent
  priors, which suggests a hybrid, but that is not tested here.
\item \textbf{Seeding.} While preparing this layer we found that the package's
  seeding helper derived its seeds from Python's builtin \texttt{hash}, which is
  salted per process for strings (PEP 456). Since every call mixes in a string
  tag, no experiment in the package was bit-reproducible from a fresh
  interpreter, contrary to what its documentation claimed. The paired
  ``common random numbers'' property that the method comparisons actually rely
  on held \emph{within} each run, so previously reported numbers stand as
  measurements; but they were not reproducible. The helper now uses a fixed
  digest and the experiments here were re-run under it. Earlier layers'
  committed outputs have not been regenerated.
\item \textbf{The state space must be represented.} The $M^2$ in every cost is
  the price of continuous messages, and it is what makes the extension beyond
  scalar-valued sites nontrivial. Nothing in Section~\ref{sec:em} requires a
  grid --- $\Xibf$ is defined by \eqref{eq:xi} for any message representation ---
  but a grid is what is implemented.
\item \textbf{One dataset, one structure.} The experiments use a single chain
  family (linear AR(1) with non-Gaussian innovations) at $n = 32$. The claim
  ``EM beats a score network'' is demonstrated in this setting, not in general;
  in particular the network's disadvantage would shrink if the prior were
  complicated enough that the chain family were itself a poor description.
\item \textbf{The baseline is a vanilla network, as specified --- and that is a
  real restriction.} The comparison uses a fully connected MLP, sweeping width
  and training budget and reporting both standard parameterizations (noise
  prediction and clean-signal prediction). It does \emph{not} use an
  architecture that encodes sequence structure. A temporal convolution or a
  U-Net would carry a locality prior of its own and should be expected to close
  part of the gap; what it would not acquire is the exactness of
  \eqref{eq:risk} or the uniformity in $t$ of Proposition~\ref{prop:sensitivity}.
  Establishing where a structured \emph{architecture} lands relative to a
  structured \emph{estimator} is the natural next comparison and is not settled
  here.
\end{enumerate}

% =====================================================================
\section{Relation to the rest of the project}

Layer 2 established that BP computes the exact score for a \emph{known} Markov
prior, and that Gaussianity buys closure of the message representation, not
exactness. Layer 4 established that for approximately Markov priors the score
correction is low rank, and that a residual network on top of a frozen
Markov-BP score is far more efficient than a direct score network. This layer
completes the sequence by removing the last thing that was handed to BP --- the
prior itself --- and shows that the same tree structure which made the score
exact also makes the \emph{learning} exact, in the sense that the E-step
involves no approximation.

Read together the three results form one argument: structural knowledge about
the data should enter the algorithm, not the training set. Layer 2 puts it in
the inference, Layer 4 puts the residual in a network, and Layer 5 puts the
estimation in a low-dimensional parameter. Each step converts a piece of what a
black-box denoiser would have to learn into something that is computed instead.

\appendix
\section{Notation}

\begin{center}
\begin{tabular}{ll}
\toprule
symbol & meaning \\
\midrule
$a \in \R^n$ & clean sequence (latent) \\
$x \in \R^n$ & noisy observation at level $t$ \\
$\alpha_t, \Delta_t$ & $e^{-t}$, $1-e^{-2t}$ \\
$K_\theta(a'\mid a)$ & transition kernel of the clean chain \\
$\ell^t_i(a_i;x_i)$ & per-site likelihood $\Normal(x_i;\alpha_t a_i, \Delta_t)$ \\
$L_i, R_i$ & forward / backward BP messages (excluding the local factor) \\
$b_i$, $b_{i-1,i}$ & single-site and pairwise posterior beliefs \\
$u_1..u_M$, $w$ & message grid and quadrature weights \\
$\Xibf \in \R^{M\times M}$ & expected pair-count matrix \eqref{eq:xi} \\
$W = \sum \Xibf$ & total edge count $N(n-1)$ \\
$Q(\theta\mid\theta')$ & expected complete-data log-likelihood \\
$J_t$ & Fisher information of $\theta$ per noisy chain \\
$m(x,t)$, $s(x,t)$ & posterior mean (denoiser) and score \\
\bottomrule
\end{tabular}
\end{center}

\section{Implementation map}

\begin{center}
\begin{tabular}{ll}
\toprule
object in this note & code \\
\midrule
E-step, $\Xibf$, evidence & \texttt{src/em.py:e\_step}, \texttt{e\_step\_multi} \\
$t\to0$ limit (clean data) & \texttt{src/em.py:clean\_statistics}, \texttt{fit\_clean} \\
EM driver, monotonicity trace & \texttt{src/em.py:fit\_em}, \texttt{EMTrace} \\
Fisher identity \eqref{eq:fisher} & \texttt{src/em.py:q\_gradient} \\
Rungs 1--4 and their M-steps & \texttt{src/kernels.py} \\
BP denoiser / DSM baseline & \texttt{src/denoiser.py} \\
Correctness, information, rates & \texttt{experiments/exp\_06\_em\_parameter\_recovery.py} \\
EM-BP vs score network & \texttt{experiments/exp\_07\_em\_vs\_score\_network.py} \\
\bottomrule
\end{tabular}
\end{center}

\end{document}
